\section{Background}

The emergence of multi-agent AI systems as a practical deployment paradigm introduces a constellation of engineering challenges that were largely theoretical a decade ago. Chief among these is the problem of \textbf{agent lifecycle management} — how autonomous agents are spawned, delegated to, monitored, and terminated within a coordinated system — and the parallel problem of \textbf{accountability and provenance}, which asks: when an agent (or chain of agents) produces a decision or action, who is responsible, and how can the chain of custody be reconstructed with confidence? These two problems are deeply intertwined. A system that spawns sub-agents without maintaining an immutable record of the parent-child relationship cannot later answer questions about responsibility, intent, or data origin. This section surveys the relevant literature across five areas: agent lifecycle patterns in multi-agent systems, accountability and provenance in distributed AI, fixed versus mutable delegation chains, hierarchical task decomposition, and the BX3 Framework as a substrate for immutable parent chains.

% --------------------------------------------------------------------
\subsection{Agent Lifecycle Management in Multi-Agent Systems}

Modern multi-agent architectures decompose complex objectives across hierarchies of specialized agents, where a parent orchestrator delegates sub-tasks to spawned agents with potentially different capabilities, contexts, and access scopes. This pattern — variants of which appear in frameworks such as LangChain's agent executor~\cite{langchain}, CrewAI's role-based agents~\cite{crewai}, and the MAKER pattern for medical diagnosis~\cite{maker-pattern} — has revealed that lifecycle management is far more than a deployment detail. It determines whether agent systems are debuggable, auditable, and safe to operate at scale.

The central lifecycle challenge is that agents are not static processes. They may be created dynamically, inherit partial context from a parent, execute asynchronously, and produce results that feed into subsequent agents. AgentSpawn (2025) formalizes this dynamic by proposing architectures where agents transfer memory to spawned children and adapt collaboration strategy based on execution feedback~\cite{agentspawn}. This work identifies that the design of the spawn mechanism — whether blocking or non-blocking, whether context is deep-copied or referenced, whether the child agent is ephemeral or persistent — has downstream consequences for both reliability and accountability.

Practical production systems have converged on a layered architecture pattern~\cite{multi-agent-architecture} comprising: (1) a \textbf{coordinator layer} responsible for task decomposition and agent spawning; (2) a \textbf{tool layer} that exposes sub-agents as composable tools to parent agents; (3) an \textbf{execution layer} that handles async result aggregation; and (4) a \textbf{policy layer} that enforces permission inheritance, rate limiting, and checkpointing. Critical lifecycle operations — cascading termination, orphan detection, cost tracking, and checkpoint/resume — are identified as production hardening requirements that most research systems neglect~\cite{production-hardening}. Without explicit lifecycle primitives, spawned agents that encounter errors or timeouts become orphaned processes whose state is unrecoverable and whose actions are untraceable.

% --------------------------------------------------------------------
\subsection{Accountability and Provenance in Distributed AI Systems}

The accountability problem in multi-agent systems is not merely one of governance philosophy — it is a technical challenge in distributed systems design. A 2025 auditability framework for LLM-based agent systems identifies five dimensions that determine whether a system can be held accountable post-deployment: action recoverability, lifecycle coverage, policy checkability, responsibility attribution, and evidence integrity~\cite{auditability-framework}. Of these, the most operationally complex is \textbf{responsibility attribution} in multi-agent chains: when Agent C produces an output that traces back through Agent B (which delegated to C) to Agent A (which was spawned by a human), the system must be able to attribute the full chain of decisions and actions to the appropriate principals.

A parallel concern is \textbf{provenance} — the tamper-evident record of what data, models, and governance controls influenced each decision. Provenance-aware AI systems propose structured audit trace collection, where provenance logs capture not just inputs and outputs but the full context of agent reasoning and delegation decisions~\cite{provenance-question}. The EU AI Act's record-keeping expectations for high-risk systems have accelerated interest in provenance architectures that can survive regulatory scrutiny without adding prohibitive overhead; prior work reports median enforcement costs of approximately 8.3ms for pre-execution mediation with tamper-evident logging, suggesting that strong accountability guarantees are achievable in production systems~\cite{auditability-framework}.

The most recent body of work addresses the specific problem of \textbf{delegation provenance}. The Human Delegation Provenance (HDP) protocol, currently an IETF Internet-Draft, proposes a lightweight cryptographic token format that encodes human authorization events and chains them across successive agent delegations as signed, append-only hops~\cite{hdp-protocol}. HDP enables offline verification of the full provenance chain using only the issuer's Ed25519 public key and a session identifier, without requiring central registries or third-party trust anchors — addressing a critical gap in OAuth 2.0 Token Exchange, JWTs, and UCAN, which all break down in multi-hop scenarios where one agent spawns another that spawns a third~\cite{multihop-delegation}. A complementary architecture from Beyer \textit{et al.} proposes human-anchored agent identity as a durable root from which explicit delegation semantics and provenance structures can be traced across ecosystems and platforms~\cite{agent-identity-architecture}.

Critically, recent work also frames \textbf{intelligent AI delegation} as an adaptive framework combining dynamic task decomposition, explicit authority transfer, continuous monitoring, and feedback-driven adjustment to handle environmental changes and unexpected failures — moving beyond brittle, heuristic delegation toward resilient, provenance-aware, auditable delegation networks suitable for agentic web deployments~\cite{intelligent-delegation}.

% --------------------------------------------------------------------
\subsection{Fixed vs. Mutable Delegation Chains}

The question of whether a delegation chain should be fixed at spawn time or allowed to mutate dynamically is a foundational design decision with consequences for both security and auditability. The dominant pattern in existing frameworks is \textbf{mutable delegation}, where parent agents may re-delegate, escalate, or transfer authority to other agents during execution. This provides flexibility but creates a significant provenance problem: if the delegation chain can change at runtime, the historical record of who authorized what becomes ambiguous, and post-hoc auditing requires reconstructing not just what happened but what the authorization state was at each decision point.

In contrast, \textbf{fixed delegation chains} — in which the parent-child relationship is immutably established at spawn time and cannot be altered — trade dynamic flexibility for strong auditability guarantees. SentinelAgent (2026) formalizes this tradeoff through its Delegation Chain Calculus (DCC), which defines seven properties for verifiable delegation in federal multi-agent AI systems, including \textit{forensic reconstructibility} (the chain can be completely reconstructed from logs) and \textit{authority narrowing} (each hop in the delegation chain can only reduce, never expand, the authority of the previous hop)~\cite{sentinelagent}. The Intent-Preserving Delegation Protocol (IPDP) operationalizes these properties through a non-LLM Delegation Authority Service that enforces pre-execution intent checking, at-execution scope enforcement, and post-execution output validation — achieving near-perfect detection with zero false positives on a 516-scenario benchmark across 10 attack categories~\cite{sentinelagent}.

The tension between fixed and mutable chains maps directly to the broader debate in distributed systems between append-only event sourcing (which favors immutability and auditability) and mutable state architectures (which favor performance and flexibility). What prior work has not fully resolved is a practical mechanism that preserves the immutability of the parent-child chain while still allowing the spawned agent to exercise dynamic judgment within its defined scope. This gap is precisely what the BX3 Framework's immutable parent chain model addresses.

% --------------------------------------------------------------------
\subsection{Hierarchical Task Decomposition in AI}

Multi-agent systems do not operate in isolation — their spawning behavior is triggered by the need to decompose complex tasks into manageable, parallelizable sub-tasks. Hierarchical Task Network (HTN) planning, originally developed in classical AI, has experienced renewed interest as a structuring principle for LLM-based agents, enabling localized replanning, parallel sub-task execution, and modular failure recovery~\cite{htn-planning}. The field has converged on a set of complementary techniques: \textbf{task decomposition strategies} that break goals into sub-tasks, \textbf{hierarchical planning structures} that organize sub-tasks into executable trees, \textbf{execution patterns} that govern plan-real-world feedback interaction, and \textbf{replanning mechanisms} that handle the inevitable divergence between plans and reality~\cite{goal-decomposition}.

A critical observation from recent work is that static decomposition is insufficient. The ADaPT framework introduces \textit{as-needed decomposition}, where a separate planner module leverages an LLM to create high-level plans only when the executor agent cannot directly handle a sub-task — recursively decomposing complex sub-tasks on an as-needed basis rather than pre-planning the entire hierarchy~\cite{adapt-decomposition}. ReAcTree extends this by using a dynamic tree structure where each agent node retrieves goal-specific experiences from episodic memory and shares information via working memory, enabling hierarchical task planning that adapts to execution feedback in real time~\cite{reactree}. The ToolTree framework applies Monte Carlo Tree Search to tool selection in LLM agents, using dual feedback from execution results and plan coherence scoring to prune the search tree during decomposition~\cite{tool-tree}.

The common thread across this literature is that effective task decomposition is not purely a planning problem — it is deeply entangled with the lifecycle and accountability of the agents that execute the decomposed sub-tasks. A decomposition strategy that spawns agents without maintaining an immutable parent chain produces a tree of actions that cannot be audited to its root. Conversely, a system with immutable parent chains but primitive decomposition logic will spawn agents inefficiently. The synthesis of these concerns — hierarchical decomposition with immutable delegation provenance — is the open problem this work addresses.

% --------------------------------------------------------------------
\subsection{The BX3 Framework as Substrate for Immutable Parent Chains}

The BX3 Framework, introduced in prior work as a universal architecture of functional roles for AI agent systems, provides the structural substrate on which immutable parent chains can be implemented~\cite{bx3-framework}. BX3 defines a role taxonomy — \textit{Orchestrator}, \textit{Executor}, \textit{Validator}, and \textit{Escalator} — that creates a deterministic hierarchy of agent responsibilities, ensuring that every agent knows its functional role relative to its parent and its children. This hierarchy is not merely organizational; it encodes enforcement semantics: each role has a defined boundary of action that its children cannot exceed, and that boundary is cryptographically verifiable through the framework's Forensic Ledger Architecture~\cite{forensic-ledger}.

Crucially, BX3's immutable parent chain property emerges from the intersection of two design decisions. First, the framework mandates that an agent may only be spawned by one parent, creating a strict tree structure (rather than a DAG) where every agent has at most one immediate predecessor. Second, the framework's Sandbox Gate Protocol enforces that spawning is a one-time, atomic event: once Agent A spawns Agent B, the relationship is sealed in the Forensic Ledger and cannot be retroactively modified. Agent B inherits its parent's role context and access scope at spawn time; it cannot request a different scope from its parent after initialization, and it cannot propagate a modified scope to its own children.

This design intentionally parallels the \textbf{append-only event sourcing} pattern in distributed systems, where state changes are recorded as immutable facts rather than mutable updates. The Forensic Ledger provides the tamper-evident, hash-chained record that enables forensic reconstruction of the full delegation chain — satisfying SentinelAgent's \textit{forensic reconstructibility} property and HDP's append-only provenance requirements within a single framework. The result is a spawning model where every agent carries its complete ancestry with it, enabling any downstream consumer — auditors, regulators, or other agents — to verify the full chain of custody without requiring access to external registries or centralized authority services.

\bibliographystyle{plain}
\bibliography{references}